%%%%%%%%%%%%
%
% $Autor: Wings $
% $Datum: 2019-03-05 08:03:15Z $
% $Pfad: TemplateSensor $
% $Version: 4250 $
% !TeX spellcheck = en_GB/de_DE
% !TeX encoding = utf8
% !TeX root = filename 
% !TeX TXS-program:bibliography = txs:///biber
%
%%%%%%%%%%%%

% Structure
\chapter{MOSFET PWM-Schaltmodul}

\section{Einleitung}

Das MOSFET PWM-Schaltmodul dient zur Steuerung elektrischer Lasten mit Hilfe eines Mikrocontrollers. 
MOSFET-Module werden in der Leistungselektronik eingesetzt, da sie hohe Ströme bei geringem Steuerstrom schalten können. Dadurch eignen sie sich zur direkten Ansteuerung durch Mikrocontroller \cite{Alles:2024}.

Durch die Ansteuerung mit Pulsweitenmodulation (PWM) kann die mittlere Leistung eines Verbrauchers geregelt werden. 
Dadurch lassen sich Motoren, Heizsysteme oder LED-Schaltungen gezielt steuern \cite{Tietze:2023}.

\section{Grundlegende Informationen}

MOSFETs (Metal-Oxide-Semiconductor Field-Effect Transistors) gehören zur Gruppe der spannungsgesteuerten Halbleiterbauelemente. 
Sie werden als elektronische Schalter oder zur Leistungsregelung eingesetzt \cite{Alles:2024}.

Ein MOSFET besitzt die drei Anschlüsse Gate, Drain und Source. 
Der Source-Anschluss dient als Bezugspotential, der Drain-Anschluss ist mit der zu schaltenden Last verbunden.
Über den Gate-Anschluss wird die Leitfähigkeit zwischen Drain und Source gesteuert. 
Liegt eine ausreichende Gate-Source-Spannung an, entsteht ein leitfähiger Kanal, durch den Strom fließen kann \cite{Tietze:2023}.

Zur Leistungsregelung wird die Pulsweitenmodulation verwendet. 
Dabei wird ein Verbraucher periodisch ein- und ausgeschaltet. 
Durch die Veränderung des Tastverhältnisses lässt sich die mittlere Ausgangsleistung beeinflussen \cite{Tietze:2023}.

Für Schaltanwendungen wird der MOSFET möglichst im ohmschen Bereich betrieben. 
In diesem Bereich verhält sich der MOSFET näherungsweise wie ein kleiner Widerstand zwischen Drain und Source \cite{Alles:2024}.

\section{Spezifisches MOSFET-Modul}

Verwendet wird ein JOY-IT MOSFET PWM-Schaltmodul. 
Das Modul dient zur Schaltung und Regelung von Gleichstromverbrauchern über ein PWM-Signal.

Auf dem Modul ist ein Leistungs-MOSFET vom Typ IRF9540N verbaut. 
Durch den geringen Durchlasswiderstand können höhere Ströme mit vergleichsweise geringer Verlustleistung geschaltet werden.

Das MOSFET-Modul wird mit einem Mikrocontroller angesteuert und eignet sich zur Steuerung von Lasten mit höherer Stromaufnahme \cite{JoyITMOSFET:2025}.

\section{Spezifikation}

Die technischen Daten wurden dem Datenblatt des Herstellers entnommen \cite{JoyITMOSFET:2025}.

\begin{table}[htbp]
	\centering
	\caption{Technische Daten des MOSFET PWM-Schaltmoduls}
	\label{tab:mosfet_specification}
	\begin{tabular}{ll}
		\hline
		\textbf{Eigenschaft} & \textbf{Wert} \\
		\hline
		Eingangsspannung & DC 5\,V bis 36\,V \\
		Steuerspannung & 3,3\,V bis 20\,V \\
		PWM-Frequenzbereich & 0\,kHz bis 20\,kHz \\
		Maximale Ausgangsleistung & 400\,W \\
		Maximaler Dauerstrom & 15\,A \\
		Betriebstemperatur & -40\,°C bis +85\,°C \\
		Abmessungen & 34\,mm $\times$ 17\,mm $\times$ 12\,mm \\
		\hline
	\end{tabular}
\end{table}

\subsection{Eigenschaften des IRF9540N}

Der verwendete IRF9540N besitzt folgende Eigenschaften:

\begin{table}[htbp]
	\centering
	\caption{Eigenschaften des verwendeten IRF9540N}
	\label{tab:mosfet_irf9540n}
	\begin{tabular}{ll}
		\hline
		\textbf{Eigenschaft} & \textbf{Beschreibung} \\
		\hline
		MOSFET-Typ & p-Kanal-Leistungs-MOSFET \\
		Funktion & elektronischer Schalter \\
		Eingangsverhalten & hoher Eingangswiderstand \\
		Leitverhalten & geringer Drain-Source-Durchlasswiderstand \\
		Anwendung & geeignet für PWM-Anwendungen \\
		\hline
	\end{tabular}
\end{table}

Der MOSFET wird auf dem Modul zur Schaltung der angeschlossenen Last verwendet.

\subsection{Anschlüsse}

Das Modul besitzt folgende Anschlüsse:

\begin{itemize}
	\item \PYTHON{SIG}
	\item \PYTHON{VCC}
	\item \PYTHON{GND}
	\item \PYTHON{VIN}
	\item \PYTHON{VOUT}
\end{itemize}

Das PWM-Steuersignal wird über den Anschluss \PYTHON{SIG} eingespeist. 
Die Versorgung der Last erfolgt über die Schraubklemmen \PYTHON{VIN} und \PYTHON{VOUT} \cite{JoyITMOSFET:2025}.

\section{Bibliothek}

\subsection{Beschreibung}

Für das MOSFET PWM-Schaltmodul wird keine zusätzliche Bibliothek benötigt. 
Die Ansteuerung erfolgt über die Standardfunktionen der Arduino-IDE.

\subsection{Funktionen}

Zur Steuerung des Moduls werden folgende Arduino-Funktionen verwendet:

\begin{itemize}
	\item \PYTHON{pinMode()}: Konfiguriert einen Pin als Ein- oder Ausgang.
	\item \PYTHON{digitalWrite()}: Setzt einen digitalen Ausgang auf \PYTHON{HIGH} oder \PYTHON{LOW}.
	\item \PYTHON{analogWrite()}: Gibt ein PWM-Signal an einem PWM-fähigen Pin aus.
	\item \PYTHON{delay()}: Erzeugt eine zeitliche Verzögerung im Programmablauf.
\end{itemize}

Mit \PYTHON{digitalWrite()} wird die Last ein- oder ausgeschaltet. 
Mit \PYTHON{analogWrite()} wird ein PWM-Signal erzeugt, wodurch sich die mittlere Ausgangsleistung regeln lässt.

\section{Beispiel}

In diesem Beispiel wird eine elektrische Last über ein PWM-Signal angesteuert.

\subsection{Versuchsaufbau}

Das PWM-Signal des Mikrocontrollers wird mit dem Eingang \PYTHON{SIG} des MOSFET-Moduls verbunden. 
Der Anschluss \PYTHON{GND} des Moduls muss mit \PYTHON{GND} des Mikrocontrollers verbunden werden, damit beide Schaltungen ein gemeinsames Bezugspotential besitzen. 
Die Last wird über die Schraubklemmen \PYTHON{VIN} und \PYTHON{VOUT} mit einer externen Spannungsversorgung verbunden. 
Die externe Versorgung muss zur angeschlossenen Last passen und den erforderlichen Strom liefern können.

Als Testlast kann zunächst eine LED mit Vorwiderstand oder eine kleine Gleichstromlampe verwendet werden. 
Eine leistungsstarke Last, beispielsweise ein Heizdraht, sollte erst nach erfolgreichem Funktionstest angeschlossen werden.

\subsection{Funktionsweise des Beispiels}

Der Mikrocontroller erzeugt mit der Funktion \PYTHON{analogWrite()} ein PWM-Signal. 
Das MOSFET-Modul schaltet dadurch den Stromfluss periodisch ein und aus. 
Durch die Veränderung des Tastverhältnisses wird die mittlere Leistung der Last beeinflusst.

\subsection{Einfacher Code}

Der folgende Code dient als einfacher Funktionstest des MOSFET PWM-Schaltmoduls. 
Dabei wird der Steuereingang des Moduls über einen digitalen Ausgang des Mikrocontrollers geschaltet.

Für den Test wird \PYTHON{SIG} mit dem im Code verwendeten Ausgangspin verbunden. 
\PYTHON{GND} des MOSFET-Moduls wird mit \PYTHON{GND} des Mikrocontrollers verbunden. 
Die Testlast wird über die Lastseite des Moduls angeschlossen. 
Die externe Versorgung wird über \PYTHON{VIN} eingespeist und die Last über \PYTHON{VOUT} geschaltet.

\begin{code}[H]
	\captionof{code}{Einfacher Funktionstest des MOSFET PWM-Schaltmoduls}
	\label{Code:MosfetSimpleSwitch}
	\ArduinoExternal{}{../../Code/MosfetSimpleSwitch/MosfetSimpleSwitch.ino}
\end{code}



\section{Kalibrierung}

Für die Ansteuerung eines Heizdrahts muss ein geeigneter PWM-Wert experimentell bestimmt werden.

Dazu wird die PWM-Stufe schrittweise erhöht, bis die gewünschte Temperatur erreicht wird. 
Dabei dürfen die zulässige Stromaufnahme des Moduls und die Belastbarkeit der Spannungsversorgung nicht überschritten werden.

\section{Einfache Anwendung}

Das MOSFET PWM-Schaltmodul kann im Projekt zur Regelung eines Gleichstromverbrauchers eingesetzt werden. 
Eine mögliche Anwendung ist die Leistungsregelung eines Heizdrahts über PWM. 
Dadurch wird die Drahttemperatur an den verwendeten Werkstoff und die gewünschte Schnittgeschwindigkeit angepasst.

Der folgende Code gibt ein PWM-Signal an das MOSFET-Modul aus. 
Der Steuereingang \PYTHON{SIG} wird dazu mit einem PWM-fähigen Ausgang des Mikrocontrollers verbunden. 
\PYTHON{GND} des Moduls und \PYTHON{GND} des Mikrocontrollers müssen verbunden sein. 
Die Lastversorgung wird an \PYTHON{VIN} angeschlossen, während der Verbraucher über \PYTHON{VOUT} geschaltet wird.

\begin{code}[H]
	\captionof{code}{PWM-Test des MOSFET PWM-Schaltmoduls}
	\label{Code:MosfetPWMTest}
	\ArduinoExternal{}{../../Code/MosfetPWMTest/MosfetPWMTest.ino}
\end{code}

\section{Tests}

\subsection{Einfacher Funktionstest}

Beim einfachen Funktionstest wird eine Testlast angeschlossen. 
Das Modul wird für mehrere Sekunden eingeschaltet und anschließend wieder ausgeschaltet.

Dabei wird überprüft, ob die Last korrekt schaltet, ob die Verdrahtung passend ausgeführt wurde und ob sich das Modul während des kurzen Tests auffällig erwärmt. 
Der Test gilt als bestanden, wenn die Last eindeutig ein- und ausgeschaltet wird und keine unzulässige Erwärmung an Modul, Leitungen oder Anschlüssen auftritt.

\subsection{Test aller Funktionen}

Beim vollständigen Test werden mehrere PWM-Werte nacheinander ausgegeben. 
Dabei wird überprüft, ob sich die Leistung der angeschlossenen Last abhängig vom Tastverhältnis verändert. 
Bei kleinen PWM-Werten muss die Last nur mit geringer Leistung betrieben werden, während bei höheren PWM-Werten eine deutlich stärkere Wirkung erkennbar sein muss.

Zusätzlich wird kontrolliert, ob das Modul bei einem digitalen Schaltsignal zuverlässig ein- und ausschaltet. 
Dazu wird die Last zunächst vollständig ausgeschaltet und anschließend für eine definierte Zeit eingeschaltet. 
Während des Tests werden die Erwärmung des MOSFET-Moduls, die Erwärmung der Anschlussleitungen und die Stromaufnahme der Last beobachtet. 

Der folgende Code testet die Heizdrahtansteuerung mit dem MOSFET PWM-Schaltmodul. 
Dabei werden ein PWM-fähiger Ausgang des Mikrocontrollers, der Steuereingang \PYTHON{SIG}, die gemeinsame Masseverbindung und die Lastseite des Moduls geprüft.

Der Heizdraht wird über die Lastseite des Moduls geschaltet. 
\PYTHON{SIG} wird mit dem verwendeten PWM-Ausgang des Mikrocontrollers verbunden. 
Zusätzlich wird \PYTHON{GND} des Moduls mit \PYTHON{GND} des Mikrocontrollers verbunden.

\begin{code}[H]
	\captionof{code}{Testprogramm für die Heizdrahtansteuerung}
	\label{Code:MosfetHotwireTest}
	\ArduinoExternal{}{../../Code/MosfetHotwireTest/MosfetHotwireTest.ino}
\end{code}

\section{Weiterführende Literatur}

\begin{itemize}
	\item Alles, Martin: \textit{Einführung in die elektronische Schaltungstechnik}. Springer Vieweg, 2024.

\end{itemize}